% =============================================================================
%  CHAPTER 4: Tables
% =============================================================================

\chapter{Tables}

\section{Basic Table (tabular)}

The \texttt{tabular} environment is the fundamental table environment in \latex{}.

\subsection{Column Types}

\begin{table}[H]
\centering
\begin{tabular}{cl}
\toprule
\textbf{Spec} & \textbf{Description} \\
\midrule
\texttt{l} & Left-aligned column \\
\texttt{c} & Centered column \\
\texttt{r} & Right-aligned column \\
\texttt{p\{width\}} & Paragraph column with fixed width \\
\texttt{|} & Vertical line \\
\texttt{@\{text\}} & Replace inter-column space with \textit{text} \\
\bottomrule
\end{tabular}
\caption{Column type specifications}
\end{table}

\subsection{Basic Example}

\begin{lstlisting}[language={[LaTeX]TeX}, caption=Basic tabular]
\begin{tabular}{|l|c|r|}
    \hline
    Left & Center & Right \\
    \hline
    A    & B      & C \\
    \hline
\end{tabular}
\end{lstlisting}

Result:
\begin{tabular}{|l|c|r|}
    \hline
    Left & Center & Right \\
    \hline
    A    & B      & C \\
    \hline
\end{tabular}

% ------------------------------------------------------------------
\section{Table Syntax}

\begin{lstlisting}[language={[LaTeX]TeX}, caption=Table syntax reference]
% &      -- column separator
% \\     -- end of row
% \hline -- horizontal line (spans full width)
% \cline{2-3} -- partial horizontal line (columns 2 to 3)

\begin{tabular}{|l|c|r|}
    \hline
    Name  & Age & Score \\
    \hline
    Alice & 25  & 95 \\
    Bob   & 30  & 87 \\
    \cline{2-3}
    \multicolumn{1}{c|}{} & \multicolumn{2}{c}{Partial line} \\
    \hline
\end{tabular}
\end{lstlisting}

% ------------------------------------------------------------------
\section{Professional Tables with booktabs}

The \texttt{booktabs} package produces publication-quality tables without vertical lines.

\begin{lstlisting}[language={[LaTeX]TeX}, caption=booktabs example]
\usepackage{booktabs}

\begin{tabular}{lcc}
    \toprule
    Name  & Age & Score \\
    \midrule
    Alice & 25  & 95 \\
    Bob   & 30  & 87 \\
    Carol & 28  & 92 \\
    \bottomrule
\end{tabular}
\end{lstlisting}

Result:
\begin{table}[H]
\centering
\begin{tabular}{lcc}
\toprule
\textbf{Name}  & \textbf{Age} & \textbf{Score} \\
\midrule
Alice & 25  & 95 \\
Bob   & 30  & 87 \\
Carol & 28  & 92 \\
\midrule
\textbf{Average} & \textbf{27.7} & \textbf{91.3} \\
\bottomrule
\end{tabular}
\caption{Student scores using booktabs}
\end{table}

\begin{table}[H]
\centering
\begin{tabular}{ll}
\toprule
\textbf{booktabs command} & \textbf{Purpose} \\
\midrule
\verb|\toprule|    & Top rule (thicker) \\
\verb|\midrule|    & Middle rule (medium) \\
\verb|\bottomrule| & Bottom rule (thicker) \\
\verb|\cmidrule{a-b}| & Partial rule (thinner) \\
\verb|\addlinespace|   & Extra vertical space \\
\bottomrule
\end{tabular}
\caption{booktabs commands}
\end{table}

% ------------------------------------------------------------------
\section{Multi-column and Multi-row}

\subsection{Multi-column (\texttt{\textbackslash multicolumn})}

\begin{lstlisting}[language={[LaTeX]TeX}, caption=Multicolumn]
\multicolumn{number_of_columns}{alignment}{text}
\end{lstlisting}

\begin{table}[H]
\centering
\begin{tabular}{llcc}
\toprule
 & \textbf{Math} & \textbf{Physics} & \textbf{Chemistry} \\
\midrule
Alice  & 95 & 88 & 91 \\
Bob    & 87 & 92 & 85 \\
\midrule
\multicolumn{2}{l}{\textit{Class Average}} & \multicolumn{2}{c}{\textbf{89.7}} \\
\bottomrule
\end{tabular}
\caption{Multicolumn example}
\end{table}

\subsection{Multi-row (multirow package)}

\begin{lstlisting}[language={[LaTeX]TeX}, caption=Multirow]
\usepackage{multirow}
\multirow{number_of_rows}{width}{text}
\end{lstlisting}

\begin{table}[H]
\centering
\begin{tabular}{llcc}
\toprule
\textbf{Category} & \textbf{Subject} & \textbf{Students} & \textbf{Avg Score} \\
\midrule
\multirow{2}{*}{Science} & Physics  & 45 & 88 \\
                          & Chemistry & 38 & 85 \\
\midrule
\multirow{2}{*}{Arts}    & History  & 32 & 91 \\
                          & Literature & 29 & 93 \\
\bottomrule
\end{tabular}
\caption{Multirow example}
\end{table}

% ------------------------------------------------------------------
\section{Fixed-width Tables (tabularx)}

\begin{lstlisting}[language={[LaTeX]TeX}, caption=tabularx]
\usepackage{tabularx}

\begin{tabularx}{\textwidth}{lX}
    \toprule
    Name & Description \\
    \midrule
    Item 1 & This text automatically wraps to fill
             the available column width. \\
    Item 2 & Another long description that wraps
             automatically in the X column. \\
    \bottomrule
\end{tabularx}
\end{lstlisting}

Result:
\begin{table}[H]
\centering
\begin{tabularx}{0.9\textwidth}{lX}
\toprule
\textbf{Command} & \textbf{Description} \\
\midrule
\verb|\LaTeX|   & Produces the \LaTeX{} logo with proper kerning and raised ``A''. \\
\verb|\TeX|     & Produces the \TeX{} logo. \\
\verb|\today|   & Produces the current date. \\
\verb|\tableofcontents| & Generates the table of contents from sectioning commands. \\
\bottomrule
\end{tabularx}
\caption{tabularx with auto-wrapping X column}
\end{table}

% ------------------------------------------------------------------
\section{Multi-page Tables (longtable)}

\begin{lstlisting}[language={[LaTeX]TeX}, caption=longtable]
\usepackage{longtable}

\begin{longtable}{lcc}
    \caption{A long table spanning multiple pages} \\
    \toprule
    Name & Value & Note \\
    \midrule
    \endfirsthead

    \toprule
    Name & Value & Note \\
    \midrule
    \endhead

    \midrule
    \multicolumn{3}{r}{\textit{Continued on next page}} \\
    \endfoot

    \bottomrule
    \endlastfoot

    Data1 & 100 & Note1 \\
    Data2 & 200 & Note2 \\
    % ... many more rows ...
\end{longtable}
\end{lstlisting}

% ------------------------------------------------------------------
\section{Table Placement (Float)}

Tables inside \texttt{table} environment are \emph{floating bodies}:

\begin{lstlisting}[language={[LaTeX]TeX}, caption=Float placement specifiers]
\begin{table}[placement]
    % placement can be:
    % h = here       (approximately at current location)
    % t = top        (top of page)
    % b = bottom     (bottom of page)
    % p = page       (on a separate float page)
    % ! = override   (relax LaTeX's internal rules)
    % H = HERE       (requires float package -- exact placement)
    \centering
    \begin{tabular}{lc}
        ...
    \end{tabular}
    \caption{Table caption}
    \label{tab:mylabel}
\end{table}
\end{lstlisting}

% ------------------------------------------------------------------
\section{Colored Tables}

\begin{lstlisting}[language={[LaTeX]TeX}, caption=Colored table rows]
\usepackage[table]{xcolor}

\rowcolor{gray!20}       % Color entire row
\cellcolor{blue!25}       % Color single cell
\columncolor{green!15}    % Color entire column (in column spec)
\end{lstlisting}

\begin{table}[H]
\centering
\begin{tabular}{lcc}
\toprule
\rowcolor{blue!15}
\textbf{Header 1} & \textbf{Header 2} & \textbf{Header 3} \\
\midrule
Row 1, Col 1 & Row 1, Col 2 & Row 1, Col 3 \\
\rowcolor{gray!10}
Row 2, Col 1 & Row 2, Col 2 & Row 2, Col 3 \\
Row 3, Col 1 & Row 3, Col 2 & Row 3, Col 3 \\
\rowcolor{gray!10}
Row 4, Col 1 & Row 4, Col 2 & Row 4, Col 3 \\
\bottomrule
\end{tabular}
\caption{Alternating row colors}
\end{table}

% ------------------------------------------------------------------
\section{Table Quick Reference}

\begin{table}[H]
\centering
\begin{tabular}{ll}
\toprule
\textbf{Symbol/Command} & \textbf{Meaning} \\
\midrule
\verb|&|               & Column separator \\
\verb|\\|              & Row terminator \\
\verb|\hline|          & Full horizontal line \\
\verb|\cline{a-b}|     & Partial horizontal line \\
\verb|\vline|          & Vertical line (inline) \\
\verb|\multicolumn{n}{spec}{text}| & Span $n$ columns \\
\verb|\multirow{n}{w}{text}| & Span $n$ rows \\
\verb|\toprule|        & booktabs top rule \\
\verb|\midrule|        & booktabs middle rule \\
\verb|\bottomrule|     & booktabs bottom rule \\
\bottomrule
\end{tabular}
\caption{Table syntax quick reference}
\end{table}
